\section{Development experiments after the implementation audit}
\label{app:newdevelopment}

This section reports new implementation checks and mechanism experiments separately from the
archived main-text study. They do not validate the historical likelihood interpretation or
establish a molecular benefit from the corrected energy loss. The runs use committed source
snapshots, fixed initial checkpoints, dedicated trace random-number streams and recorded failures.
The development evidence and per-seed values are generated from the raw run records by
\texttt{scripts/research/render\_evidence.py}. The code currently passes 88 tests, including an
actual FlowMol endpoint head, full trajectory gradients, independent replicas and conditional FM.

\subsection{A conditional sampler that matches the density definition}

Let $m$ denote the atom count, ordered atom types and charge inputs. We separate a frozen
composition prior $\pi(m)$ from a position flow $q_{\theta,T}(x\mid m)$:
\[
 m\sim\pi,\qquad X_0\sim\mathcal N_{\mathrm{COM}}(0,I),\qquad
 \dot X_t=P v_\theta(PX_t,t;m),\quad 0\le t\le T.
\]
Here $P$ removes the unweighted centroid. There is no history-dependent self-conditioning,
prior alignment or steric retraction. A frozen original generator can supply $m$, discarding
its coordinates; the position checkpoint must not replace that composition checkpoint, since
its shared backbone updates also change the discrete logits. The conditional geometry density
is specified, while the implicit composition prior is not claimed to have a tractable density.

The initial position experiment uses $T=0.8$. To train the data to this time, rather than merely
stopping a time-one flow early, define $\tau(t)=\alpha(t)/\alpha(T)$ and
\[
 X_t=X_0+\tau(t)(X_1-X_0),\qquad
 D^*=X_0+\frac{X_1-X_0}{\alpha(T)}.
\]
The existing endpoint-to-velocity conversion then gives
\[
 \frac{\alpha'(t)}{1-\alpha(t)}(D^*-X_t)
   =\frac{\alpha'(t)}{\alpha(T)}(X_1-X_0).
\]
Endpoint regression therefore defines a positive time-weighting of the appropriate conditional
FM loss. This position objective replaces the joint CTMC objective during geometry fine-tuning.
The initial experiment sets the force and energy weights to zero; OMol25, 83 element types,
bond-free training and the 200-atom model limit remain unchanged.

Calling \texttt{eval()} alone is insufficient to disable the original network's special
self-conditioning branch exactly at $t=0$. The corrected context also disables and restores
that flag. Interior midpoint calculations are unchanged; adaptive solvers that evaluate the
boundary require this explicit change. It fixes the field definition but does not establish
numerical convergence of the trained field.

\subsection{Bias from squaring a stochastic density}

For $K$ geometries within a parent, let $C=I-\mathbf1\mathbf1^\top/K$ and
$r=\ell_h+\beta E$, where $\ell_h$ is the exact-trace value along a specified deterministic
numerical trajectory. If $\widehat r=r+\epsilon$ has conditionally zero-mean trace noise and
covariance $\Sigma$, then the standard quadratic-form identity gives
\[
 \mathbb E\,\frac{\|C\widehat r\|^2}{K}
 =\frac{\|Cr\|^2}{K}+\frac{\operatorname{tr}(C\Sigma)}{K}.
\]
The extra term generally depends on model parameters. For two conditionally independent trace
replicas along the same trajectory,
\[
 L_\times=\frac{(C\widehat r_1)^\top(C\widehat r_2)}{K},\qquad
 \mathbb E L_\times=\frac{\|Cr\|^2}{K}.
\]
Against the equal-budget squared replica mean $L_{\mathrm{mean}}$,
\[
 L_{\mathrm{mean}}-L_\times
 =\frac{\|C(\widehat r_1-\widehat r_2)\|^2}{4K}\ge0.
\]
These are standard independent-product identities, not claimed new theorems. Their gradient
interpretation additionally requires full trajectory and prior differentiation and an exchange
of differentiation and expectation. Positivity clipping or outcome-dependent skipping changes
the objective. The product can be negative and has substantial optimization variance.
Neither replica products nor common probes correct ODE integration error.

We include common probes across sibling geometries as a control, while keeping the two replicas
independent. Common probes can cancel noise for an affine field, but do not generally cancel
state-dependent Jacobian contributions. Gaussian probes control Rademacher probes' orientation
dependence; they do not remove squared-loss bias. Off-policy log-dispersion regularisation itself
is already covered by LDR \citep{ldr2026}. Molecular usefulness and computational advantages
remain necessary beyond the estimator identities.

\subsection{Two exactly solvable mechanism experiments}

The first experiment uses a two-dimensional Gaussian CNF with
$A=a\left[\begin{smallmatrix}0&1\\1&0\end{smallmatrix}\right]$ and target $a=0.6$.
It isolates trace-noise bias from numerical integration. Over five optimization seeds, the
squared two-probe residual attains mean exact KL $0.74756$, while the independent product
attains $0.00138$. Exact traces, common probes and the rotated diagonal squared-loss control
all attain zero to numerical precision. Thus this example does not demonstrate superiority
over common probes.

The second experiment uses the nonlinear shear
\[
 v_a(x)=(a\sin(3x_2),0),\quad
 \Phi_a^{-1}(y)=(y_1-a\sin(3y_2),y_2).
\]
Its exact divergence is zero, but its stochastic trace is
$3a\cos(3x_2)\xi_1\xi_2$. Common probes therefore leave state-dependent differences between
geometries. Independent off-policy examples have standard deviation $0.35$ in each coordinate;
all arms use groups of 64, two probes, 3,000 Adam steps and the same five data seeds. The target
is $a=0.6$. Ground-truth KL is
$\tfrac14(a-0.6)^2(1-e^{-18})$. We report every arm in \Cref{tab:sheartoy}; the noisier product
variants are not omitted.

\begin{table}[ht]
\centering
\caption{Nonlinear shear mechanism test. Mean exact KL over five optimization seeds; no molecular
claim. All numerical values are from completed runs.}
\label{tab:sheartoy}
\begin{tabular}{lr}
\toprule
Estimator and probe policy & Mean exact KL $\downarrow$\\
\midrule
Exact trace & $<10^{-30}$\\
Squared, independent Rademacher & 0.08547\\
Product, independent Rademacher & 0.00259\\
Squared, common Rademacher & 0.08141\\
Product, common Rademacher & 0.03745\\
Squared, independent Gaussian & 0.08732\\
Product, independent Gaussian & 0.01582\\
Squared, common Gaussian & 0.08211\\
\bottomrule
\end{tabular}
\end{table}

These constructed examples establish a mechanism and illustrate variance tradeoffs. They do
not predict the magnitude or sign of an effect in molecular training.

\subsection{Molecular runtime, split audit and preliminary sampling}

Four equal trace-budget controls---squared versus product, crossed with independent versus
common probes---each complete 20 real optimizer steps with finite weights. The two
independent-probe arms each skip one energy outlier, and the common-probe arms skip none.
Their raw objective means are not compared as performance endpoints. A separate eight-parent
numerical panel at $T=0.95$ keeps geometries and probes fixed across 16, 32, 64 and 128 steps.
Two parents still have maximum centred mean log-density changes of 90.36 and 66.18 nats between
64 and 128 steps. Those density values are not certified as converged. These evaluation probes
are fixed over time, whereas training resamples probes over integration steps; their noise
magnitudes are not interchangeable.

A full scan of 3,902,107 archived training structures finds 6,527 validation structures whose
elemental composition is absent from the training set. Of these, 59 contain at most 12 atoms.
This excludes an exact chemical-identity match with the training set, but it is still an old
validation development pool, not a new blind test set. The historical test file duplicates
validation. Original source identities and spin metadata remain unavailable.

One conditional position FM run completes 1,000 optimizer steps, processing distinct random
OMol25 training structures, in 104.2 seconds of its training loop with peak allocated memory
5.06 GiB. We select eight development parents with a fixed seed, excluding charge values at
legacy clipping boundaries. For each parent, four independent Gaussian coordinate priors are
shared between the warm checkpoint and the fine-tuned checkpoint. Both use the same
128-step clamped sampler at $T=0.8$. The warm checkpoint has not been trained for this new
conditional path, so this is a warm-start diagnostic, not a fair comparison with its original
joint generator.

Independent GFN2-xTB optimization converges for 29/32 warm samples and 30/32 fine-tuned samples.
Successful samples have median strain 13.09 and 5.01 eV, respectively; these medians condition on
different success sets. Pairing the same priors and ranking a nonconverged result below a
converged one gives 28 improvements, three deteriorations and one pair in which both fail.
The eight reference structures all converge and have median strain 1.07 eV. This is one
training seed, not a replicated energy-training result. Original spin is unknown; both
single-point and relaxation calculations explicitly use minimum electron-parity spin and the
stored charge. The result is under that declared computational convention.

Sampling also remains solver-sensitive: the median/max coordinate RMS change from 64 to 128
steps is 0.0093/0.0902 \AA{} for the warm checkpoint and 0.0089/0.1585 \AA{} after fine-tuning.
Further training, tighter numerical checks, multiple seeds, fair conditional baselines and
independent data are required before an energy-regularisation or Boltzmann-sampling claim.
Strain alone does not establish valid connectivity, diversity or equilibrium populations.
